\chapter{Introduction}

High-speed serial links suffer frequency-dependent channel loss that closes the
received eye and reduces the decision margin. Receiver equalization restores
useful high-frequency content before sampling. For a CTLE, the relevant design
variables interact: greater peaking may improve the eye after a lossy channel,
but can also increase noise, reduce headroom, or create an undesirable response.
Manual exploration becomes costly when every candidate must pass several SPICE
analyses and multiple operating conditions.

Nebula investigates whether a structured optimization loop can automate this
exploration. A candidate generator proposes circuit and DFE parameters; the
simulation layer converts the proposal into parameterized ngspice benches; and
the analysis layer returns electrical metrics, failures, and provenance. The
framework targets a \SI{5}{Gb/s} signaling context corresponding to PCIe Gen~2,
but it does not claim protocol or compliance certification \citep{pcieOverview}.

\begin{quote}
\textbf{Naming convention.} The product developed by the team is named
\emph{Nebula}. The competition is also named the \emph{Astera Labs Nebula
Competition}. To prevent ambiguity, this report uses ``Nebula'' for the project
product and writes ``Nebula Competition'' whenever it refers to the competition.
\end{quote}

The project source is available at
\href{\projecturl}{\texttt{github.com/dwiti-25/nebula}}. The repository is the
authoritative source for implementation behavior; this report distinguishes
implemented functionality from historical experimental evidence.

\section{Contributions}

The principal engineering contributions are:
\begin{enumerate}
  \item a parameterized SKY130 CTLE and separate DC, AC, transient, noise, and
        HD3 ngspice benches;
  \item a four-port Touchstone channel path with explicit differential port
        mapping and validation;
  \item a behavioral one-tap DFE with phase selection, eye-opening, margin, and
        finite-pattern error metrics;
  \item a staged evaluator with early rejection, caching, structured failures,
        and provenance fingerprints;
  \item PPO, Random Search, and CEM candidate generation under a shared
        evaluation framework; and
  \item command-line, local web, and natural-language interfaces for training,
        inference, qualification, run-graph inspection, and artifact export.
\end{enumerate}

\section{Scope}

Nebula evaluates a Stage~1 schematic model. Its transistor-level component is
the CTLE. The channel is represented by measured or synthetic S-parameters, and
the DFE is evaluated in Python after the CTLE transient simulation. The system
does not contain a transistor-level sampler, clock-data recovery, physical DFE,
layout parasitics, package model, or protocol checker. Consequently, successful
evaluation means that the executed project-defined checks passed; it does not
establish production readiness or standards compliance.
